\documentclass[11pt,a4paper]{article}

% ============================================================================
% ZIPA-CLI User Manual
% Build:  make            (uses latexmk)   ->  manual.pdf
%   or:   pdflatex manual.tex  (run twice for the table of contents)
% ============================================================================

\usepackage[utf8]{inputenc}
\usepackage[T1]{fontenc}
\usepackage{lmodern}
\usepackage[margin=1in]{geometry}
\usepackage{hyperref}
\usepackage{xcolor}
\usepackage{listings}
\usepackage{booktabs}
\usepackage{longtable}
\usepackage{enumitem}
\usepackage{fancyhdr}
\usepackage{titlesec}

\hypersetup{
  colorlinks=true,
  linkcolor=blue!55!black,
  urlcolor=blue!55!black,
  citecolor=blue!55!black,
  pdftitle={ZIPA-CLI User Manual},
  pdfauthor={zipa-cli}
}

\definecolor{codebg}{rgb}{0.96,0.96,0.96}
\definecolor{codekw}{rgb}{0.10,0.20,0.55}
\definecolor{codecm}{rgb}{0.30,0.45,0.30}

\lstdefinestyle{shell}{
  backgroundcolor=\color{codebg},
  basicstyle=\ttfamily\small,
  keywordstyle=\color{codekw}\bfseries,
  commentstyle=\color{codecm}\itshape,
  breaklines=true,
  columns=fullflexible,
  frame=single,
  rulecolor=\color{gray!40},
  showstringspaces=false,
  xleftmargin=6pt,
  morekeywords={zipa-cli,decode,transcribe,models,compare,pip,python}
}
\lstset{style=shell}

\newcommand{\cli}{\texttt{zipa-cli}}
\newcommand{\code}[1]{\texttt{#1}}

\pagestyle{fancy}
\fancyhf{}
\rhead{ZIPA-CLI User Manual}
\lhead{\leftmark}
\cfoot{\thepage}

\titleformat{\section}{\Large\bfseries}{\thesection}{0.6em}{}

\begin{document}

% ----------------------------------------------------------------------------
\begin{titlepage}
\centering
\vspace*{3cm}
{\Huge\bfseries ZIPA-CLI\par}
\vspace{0.6cm}
{\Large Batch Phonetic Decoding with ZIPA Zipformer Models\par}
\vspace{0.4cm}
{\large User Manual\par}
\vfill
{\large ONNX \& PyTorch backends \quad$\cdot$\quad multi-format I/O \quad$\cdot$\quad
timestamped alignment \quad$\cdot$\quad model comparison \quad$\cdot$\quad web viewer\par}
\vspace{2cm}
{\itshape Generated for the \cli{} package.\par}
\end{titlepage}

\tableofcontents
\newpage

% ============================================================================
\section{Introduction}
\cli{} is a command-line tool for running \emph{greedy phonetic decoding} with the
ZIPA family of zipformer phone-recognition models
(\href{https://aclanthology.org/2025.acl-long.961/}{ACL 2025}). It wraps the model
checkpoints in a single streaming pipeline so you can transcribe one file or millions
without reloading the model, across many input formats, with either a
minimal-dependency ONNX backend or the full PyTorch backend.

\paragraph{Capabilities at a glance.}
\begin{itemize}[nosep]
  \item Inputs: single file, file list, directory, CommonVoice-style TSV, STM
        segments, HuggingFace dataset, lhotse/icefall manifests, and lhotse shar.
  \item Backends: \textbf{ONNX} (default, minimal deps) and \textbf{PyTorch}
        (icefall/k2, GPU).
  \item Outputs: \code{tsv}, \code{jsonl}, lhotse \code{manifest}, \code{recogs},
        \code{stm}, \code{ctm}, and \code{align-json}.
  \item Timestamped alignment, a two-model comparison mode (match/sub/ins/del, PER,
        PFER), a model registry with HuggingFace auto-download, and a client-side web
        viewer for phoneticians.
\end{itemize}

% ============================================================================
\section{Installation}

\subsection{Minimal (ONNX backend --- recommended)}
\begin{lstlisting}
pip install -e ".[onnx]"
\end{lstlisting}
Pulls \code{onnxruntime soundfile librosa lhotse huggingface\_hub sentencepiece}. No
\code{torch}/\code{icefall}/\code{k2} required.

\subsection{Full (PyTorch backend --- large-scale GPU)}
\begin{lstlisting}
pip install -e ".[torch]"
# then install icefall + k2 matching your torch/cuda versions:
#   https://icefall.readthedocs.io/en/latest/installation/index.html
\end{lstlisting}

\subsection{Optional extras}
\begin{lstlisting}
pip install -e ".[hf]"        # HuggingFace dataset input
pip install -e ".[analysis]"  # panphon PFER in compare mode
pip install -e ".[all]"       # everything
\end{lstlisting}

\subsection{Configuration}
\begin{longtable}{@{}lll@{}}
\toprule
\textbf{Variable / flag} & \textbf{Meaning} & \textbf{Default} \\
\midrule
\endhead
\code{ZIPA\_CLI\_CACHE} & model download cache & \code{\textasciitilde/.cache/zipa-cli} \\
\code{ZIPA\_REPO} / \code{--zipa-repo} & path to the ZIPA training repo & auto-detected \\
\code{--tokens} & ONNX \code{tokens.txt} & bundled \code{ipa\_simplified/tokens.txt} \\
\code{--bpe-model} & PyTorch sentencepiece model & bundled \code{unigram\_127.model} \\
\bottomrule
\end{longtable}

% ============================================================================
\section{Models}
\cli{} knows all 12 released averaged checkpoints by short tag.

\begin{lstlisting}
zipa-cli models list
zipa-cli models info  zipa-cr-l-500k
zipa-cli models download zipa-cr-s-300k --backend onnx --precision fp16
\end{lstlisting}

\begin{longtable}{@{}lll@{}}
\toprule
\textbf{Tag pattern} & \textbf{Architecture} & \textbf{Example} \\
\midrule
\endhead
\code{zipa-t-\{s,l\}-\{300k,500k\}}    & transducer        & \code{zipa-t-l-500k} \\
\code{zipa-cr-\{s,l\}-\{300k,500k\}}   & CR-CTC (ctc)      & \code{zipa-cr-s-300k} \\
\code{zipa-cr-ns-\{s,l\}-\{700k,800k\}}& CR-CTC Ns         & \code{zipa-cr-ns-l-800k} \\
\code{zipa-cr-ns-nd-\{s,l\}-\dots}     & CR-CTC Ns, no diacritics & \code{zipa-cr-ns-nd-l-780k} \\
\bottomrule
\end{longtable}

Anywhere a model is required you may pass a \textbf{tag} (auto-downloaded and cached)
or a \textbf{local path} (an \code{.onnx} file or transducer directory, or a
\code{.pth} checkpoint).

% ============================================================================
\section{Quick start}
\begin{lstlisting}
# Single file (prints phones)
zipa-cli transcribe sample.wav --model zipa-cr-s-300k

# A directory of audio on GPU, dynamic batching, to a TSV
zipa-cli decode --input audio_dir/ --input-type dir \
    --model zipa-cr-l-500k --backend onnx \
    --max-duration 600 -o transcripts.tsv

# Compare two models' outputs
zipa-cli decode --input audio_dir/ --model zipa-cr-s-300k -o small.tsv
zipa-cli decode --input audio_dir/ --model zipa-cr-l-500k -o large.tsv
zipa-cli compare --a small.tsv --b large.tsv
\end{lstlisting}

% ============================================================================
\section{Input sources}
Select with \code{--input-type}; \code{auto} (the default) infers from the path/flags.

\begin{longtable}{@{}lll@{}}
\toprule
\textbf{Type} & \textbf{What it is} & \textbf{Key flags} \\
\midrule
\endhead
\code{file} & one audio file & --- \\
\code{list} & text file of audio paths & --- \\
\code{dir} & directory searched recursively & --- \\
\code{tsv} & id/path (+ ref) columns, CommonVoice-style & \code{--id-col --path-col --ref-col} \\
\code{stm} & \code{file spk chan start end [text]} segments & \code{--audio-dir} \\
\code{hf} & a HuggingFace dataset & \code{--hf-dataset --hf-split --hf-config} \\
\code{manifest} & lhotse/icefall manifests & \code{--cuts} \emph{or} \code{--recordings --supervisions} \\
\code{shar} & directory of lhotse shar shards & --- \\
\bottomrule
\end{longtable}

\subsection{Examples}
\begin{lstlisting}
# CommonVoice TSV (segid -> clip)
zipa-cli decode --input cv/test.tsv --input-type tsv \
    --id-col path --path-col path --ref-col sentence \
    --model zipa-cr-l-500k -o cv.tsv

# STM segments + an audio directory
zipa-cli decode --input show.stm --input-type stm --audio-dir audio/ \
    --model zipa-cr-l-500k --output-format stm -o show.hyp.stm

# HuggingFace dataset (streaming)
zipa-cli decode --input-type hf --hf-dataset mozilla-foundation/common_voice_17_0 \
    --hf-config en --hf-split test --audio-column audio \
    --model zipa-cr-l-500k -o cv17_en.jsonl --output-format jsonl

# icefall manifest (recordings + supervisions)
zipa-cli decode --input-type manifest \
    --recordings data/recordings.jsonl.gz --supervisions data/supervisions.jsonl.gz \
    --model zipa-cr-l-500k -o hyp.manifest.jsonl --output-format manifest

# Precomputed fbank cuts (feature extraction skipped automatically)
zipa-cli decode --input-type manifest --cuts data/cuts.jsonl.gz \
    --model zipa-cr-l-500k -o hyp.tsv

# lhotse shar shards (cuts.*.jsonl.gz + recording.*.tar)
zipa-cli decode --input data-shar/ --input-type shar \
    --model zipa-cr-l-500k -o hyp.tsv
\end{lstlisting}

% ============================================================================
\section{Output formats}
\begin{longtable}{@{}ll@{}}
\toprule
\textbf{Format} & \textbf{Description} \\
\midrule
\endhead
\code{tsv} (default) & \code{id<TAB>p h o n e s} \\
\code{jsonl} & rich record per utterance (add \code{--timestamps} for an \code{alignment} field) \\
\code{manifest} & lhotse \code{SupervisionSet}; round-trips into k2/lhotse \\
\code{recogs} & \code{hyp=/ref=} lines parsed by \code{zipa/scripts/evaluate.py} \\
\code{stm} & \code{file spk chan start end phones} \\
\code{ctm} & NIST CTM: \code{recording channel start dur phone} (timed) \\
\code{align-json} & one JSON line per utterance with timed phones; feeds the web viewer \\
\bottomrule
\end{longtable}

When the input carries a reference (TSV \code{--ref-col}, STM text, manifest
supervisions) it is preserved into \code{jsonl}/\code{recogs}/\code{align-json} for
evaluation. \code{--skip-existing} resumes a run by appending only new ids.

% ============================================================================
\section{Batching, device, and precision}
\begin{itemize}[nosep]
  \item \code{--batch-size N} --- fixed utterances per batch, \emph{or}
  \item \code{--max-duration S} --- dynamic, lhotse-style: cap pooled audio seconds
        per batch (utterances length-sorted within a buffer to minimise padding).
  \item \code{--num-workers N} --- audio/feature loading workers.
  \item \code{--device auto|cpu|cuda|cuda:N} --- PyTorch device; ONNX uses the CUDA
        execution provider automatically when available.
  \item \code{--precision fp32|fp16|int8} --- ONNX precision to download/use.
\end{itemize}

% ============================================================================
\section{Timestamps and alignment}
\code{ctm} and \code{align-json} always emit per-phone \code{start}/\code{end} times;
\code{--timestamps} additionally attaches an \code{alignment} array to \code{jsonl}.

\begin{lstlisting}
zipa-cli decode --input utt.wav --model zipa-cr-l-500k \
    --output-format ctm -o utt.ctm

zipa-cli decode --input audio_dir/ --input-type dir --model zipa-cr-l-500k \
    --output-format align-json -o aligned.json
\end{lstlisting}

\begin{quote}\itshape
\textbf{Caution.} Timestamped decoding currently requires the ONNX backend, and ONNX
timings are approximate --- see Section~\ref{sec:caution}.
\end{quote}

% ============================================================================
\section{Comparison / analysis mode}
\begin{lstlisting}
zipa-cli compare --a modelA.tsv --b modelB.tsv \
    [--format txt|jsonl|csv] [--top-k 20] [-o report.txt]
\end{lstlisting}
Treats \textbf{A as reference} and \textbf{B as hypothesis}, joins by id, and reports:
corpus and per-utterance \emph{Match / Substitution / Insertion / Deletion} counts,
\textbf{PER} $=(S+I+D)/N_{\text{ref}}$, \textbf{PFER} (panphon articulatory feature
edit distance; install \code{[analysis]}), the most frequent substitution pairs and
inserted/deleted phones, and the most-disagreeing utterances. Inputs may be
\code{tsv} or \code{jsonl} transcripts from \code{decode}.

% ============================================================================
\section{Web viewer for phoneticians}
A zero-install, client-side viewer lives in \code{web/}. Produce \code{align-json}
for each model, then serve the folder:
\begin{lstlisting}
zipa-cli decode --input utt.wav --model zipa-cr-l-500k --output-format align-json -o large.json
zipa-cli decode --input utt.wav --model zipa-cr-s-300k --output-format align-json -o small.json
cd web && python -m http.server 8000   # browse http://localhost:8000
\end{lstlisting}
Load the audio file and both JSON files; entries sharing an utterance id are stacked
as one phone tier per model, positioned by timecode, over a waveform/spectrogram view.

% ============================================================================
\section{Caution: the ONNX length heuristic}\label{sec:caution}
The two backends compute the number of decodable output frames differently,
\emph{by design}.

\begin{longtable}{@{}lll@{}}
\toprule
\textbf{Backend} & \textbf{Output-length source} & \textbf{Time stride / frame} \\
\midrule
\endhead
ONNX CTC        & \code{feat\_lens // 2}, clipped to session output & $\approx 0.02$\,s \\
ONNX transducer & \code{feat\_lens // 4}, clipped                   & $\approx 0.04$\,s \\
PyTorch (both)  & exact \code{encoder\_out\_lens}                   & exact \\
\bottomrule
\end{longtable}

Why \code{//2} for CTC when \code{subsampling\_factor} is 4? The CR-CTC head emits at
roughly half the encoder stride, so its log-probs come out at $\sim$2$\times$ the
encoder frame rate. The heuristic (carried verbatim from
\code{inference/batch\_inference.py}) approximates only the \emph{trailing} valid
length; it does not change the bulk of the emitted tokens but can clip or extend the
last few frames and it sets the time stride used to place phones on the timeline.

\paragraph{Recommendation.} On your first real ONNX CTC run, spot-check the
transcript and any \code{--timestamps} alignment against a known reference (or against
the PyTorch backend, which is exact). For precise timings prefer \code{--backend torch}.

% ============================================================================
\section{Faithfulness to the reference scripts}
\begin{itemize}[nosep]
  \item Features: 80-dim lhotse kaldi fbank (\code{dither=0.0, snip\_edges=False}),
        mono, 16\,kHz --- identical to \code{inference/utils.py}.
  \item Decoding is greedy everywhere (blank id 0; transducer context size 2, $\le$3
        symbols/frame).
  \item Tokenizers bound per backend: \code{tokens.txt} (ONNX) vs
        \code{unigram\_127.model} (PyTorch).
  \item Architecture/size come from the registry, not from string-matching the path.
\end{itemize}

% ============================================================================
\section{Troubleshooting}
\begin{longtable}{@{}p{0.42\textwidth}p{0.52\textwidth}@{}}
\toprule
\textbf{Symptom} & \textbf{Fix} \\
\midrule
\endhead
\emph{Could not find tokens.txt} & Pass \code{--tokens} or set \code{--zipa-repo}/\code{\$ZIPA\_REPO}. \\
\emph{No module named icefall} & PyTorch backend needs icefall+k2 installed separately. \\
\emph{No transducer encoder/decoder/joiner ONNX files} & Wrong \code{--precision} or path; check \code{models info <tag>}. \\
ONNX timings look off & Expected approximation; validate vs PyTorch (Section~\ref{sec:caution}). \\
Web imports blocked & Serve over HTTP (\code{python -m http.server}); \code{file://} blocks ES modules. \\
\bottomrule
\end{longtable}

\end{document}
